\scanpage{54}
\mspage{47}

Pour les \deleted{équivalences (2)} inclusions (2), on utilise le corollaire \inserted{de la prop. 2} pour les couples $\Phi=\mathrm{TC}(I),\Psi'=F(I')$, resp. $\Phi=C(I)$; $\Psi'=\unsure{E|F|\mathrm{TF}}(I')$, et on applique (1), qui donne $u^*(\Phi)\leftarrow\Psi'$, d'où $\Phi\longleftrightarrow u_*(\Psi')$, ce qui n'est autre que (2). Pour les équivalences (3), on utilise la proposition 2, en faisant $\Psi=F(I),\Phi'=\mathrm{TC}(I')$ resp\unsure{t}. $\Psi=\mathrm{TF}(I),\Phi'=C(I')$.

\uline{Corollaire 1} (de la proposition 1) : Supposons $u:I'\unsure{-|\to}I$ telle que pour tout $i\in\mathrm{Ob}\,I$, la catégorie $i\backslash I'$ soit discrète\textsuperscript{*} \deleted{(p. ex. \unsure{que $u$ soit} \unsure{isom.} local)} (p. ex. $I'$ discrète). Alors on a

% slanted

\margin{Plus généralement, il suffit que que \unsure{la|les|toute} \unsure{composante|composantes} \unsure{connexe|connexes} de $i\backslash I'$ ait \inserted{$e_\alpha$} un objet initial, de sorte que \[ \underset{i\backslash I'}{\varinjlim}\,X'(i')\simeq\coprod_{\alpha\in\pi_0(i\backslash I')}X'(e_\alpha). \] C'est le cas notamment si $I'\to I$ est un isom. local, p. ex. une \unsure{immers. ouverte}.}

\[ \left\{\begin{aligned}
u_!\deleted{\unsure{(}}C(I')&\subset\deleted{\unsure{T}}C(I)\unsure{)},&u_!(\mathrm{TC}(I'))&\subset\mathrm{TC}(I)\\
u^*(F(I))&\subset F(I'),&u^*(\mathrm{TF}(I))&\subset\mathrm{TF}(I')
\end{aligned}\right. \]

En effet, les relations de la deuxième ligne équivalent (à cause de (3)) à celles de la première. Il faut prouver que si $f$ dans $\mathrm{Fl}\,M(I')$ est dans $C(I')$ resp. $\mathrm{TC}(I')$, $u_!(f)$ est dans $C(I)$ resp. $\mathrm{TC}(\unsure{I|II})$, ce qui signifie aussi que $\alpha_s^*(u_!(f))$ est dans $C$ resp. $\mathrm{TC}$, pour $\forall\,s\in\mathrm{Ob}\,I$. Or le foncteur

\[ \alpha_s^*u_!:\deleted{\unsure{A}}M(I')\longrightarrow M \]

s'explicite par la formule bien connue

\[ \alpha_s^*(u_!(X'))\overset{\text{déf}}{=}u_!(X')(s)\simeq\varinjlim_{i'\in s\backslash I'}X'(i').\quad\text{Comme} \]
% archive sheet number: 53
